\section{讨论}

\subsection{原始模型的能力边界}

本次统一审计显示，原始 Frankenstein-v28 的甲基化头并非完全失效：修正真值后，已知序列与端到端 AUC 仍分别达到 0.900 和 0.908，说明预测分数对阳性、阴性位点具有排序信息。然而，固定在项目预设阈值 $>0.6$ 时，Recall 只有 46.36\% 和 42.53\%，意味着超过一半真实甲基位点没有被保留。与此同时，天然化 RAA 只有 16.08\%，所以“甲基分类不错”和“完整扩展氨基酸序列恢复好”不能互相替代。

六个非 3AV 复合物给出了更接近实际筛选流程的答案。544 条温度 0.5 原始序列中，476 条通过甲基保留门；联合 RMSD $<5$~\AA 的最终端到端产率为 18.57\%，$<3$~\AA 为 2.94\%。如果只对 476 条保留样本报告 21.22\% 和 3.36\%，会忽略甲基门带来的损失，因此本文同时固定两个分母。

结构误差的主要来源也很清楚：476 条中 469 条全局复合物 CA RMSD $<3$~\AA、全部 476 条 $<5$~\AA，但联合通过数只有 16 和 101。换言之，瓶颈不是整体受体框架，而是同一全局坐标系内的环肽结合姿态。best85 中“肽自拟合约 2.65~\AA、受体框架拟合后约 29.93~\AA”的差距从另一队列支持了同一解释，但因选择方式和 RMSD 定义不同，不能把两个队列合并计算。

\subsection{不同靶点需要不同诊断}

六个靶点没有表现出统一失败模式。3WNE 的甲基保留率只有 27/84（32.14\%），且保留后环肽 RMSD 中位数达到 29.00~\AA，甲基门和结合姿态都是瓶颈。3ZGC 的保留率为 49/60（81.67\%），联合 $<5$~\AA 的原始产率达到 38.33\%，但其全局 RMSD 中位数 2.34~\AA 高于其他靶点，提示整体结构拟合仍需关注。1SFI、3P8F、4K1E 和 4KEL 全部通过甲基保留门，但联合 $<5$~\AA 的原始产率仍只有 10--24\%，说明仅提高甲基标记概率并不能解决姿态问题。

因此，后续模型改进不宜只优化一个总体 Accuracy。更合理的错误分解是：基础残基类别、甲基化判断、完整序列有效率、受体框架质量、环肽内部形状和环肽结合姿态分别追踪，并按靶点报告。只有这样才能判断修改真正修复了哪一步。

\subsection{数据审计为何改变结论}

旧单体标签把 62 个 Ser 错列为甲基阳性，旧评价脚本又把批次 padding 的零初始化位置当作真实阴性。前一个问题改变真值，后一个问题使结果依赖 batch size。两者都会产生表面上更好、却不可复现的 Accuracy 或 F1。本文没有改变检查点和位点预测概率，只修正评价标签并限制到 1505 个真实位点，因此主结果应视为对同一模型的更干净估计，而不是一次新的训练结果。

数据来源审计还改变了单体集合的性质判断。规划截图提出了 ``Baker 33'' 和 ``Rosetta 400+300''，但逐文件复核显示 751/751 个原始 PDB 均声明为 Rosetta 理论模型，其中406个来自 Rosetta-2023、345个来自 Rosetta-2025；没有找到可核验的 Baker 33 实验结构子集。因此，截图应被视为早期规划，而不能覆盖冻结数据本身。当前仍缺少的是 checkpoint 的完整训练 manifest 与同源聚类级泄漏审计，而不是这些 PDB 的 Rosetta 理论来源。

\subsection{限制与尚未完成的论文指标}

本稿的主要限制如下：

\begin{enumerate}
  \item 151 个单体没有与原始 V28 一一对应的预测结构结果，因此不能报告单体 CA/backbone/all-atom RMSD、pLDDT、能量或通透性。历史 164 条结构表属于 V12 工作流，不能移植到 V28。
  \item 六复合物主分析有 CA 层面的全局与环肽 RMSD，但没有统一可用的 backbone/all-atom 甲基化 RMSD、结合位点恢复、同靶点完整多样性或原生对照能量。
  \item best85 使用真实天然序列恢复率挑选，存在 oracle 选择偏倚；其结果不能代表盲生成或 100 条/靶点的未来批次。
  \item 置信度、通透性与 Rosetta 分数均为计算量，尚无实验结构、实验通透性、亲和力或稳定性验证。
  \item 主要比例来自 6 个靶点，靶点间异质性明显；逐序列 Wilson 区间不包含靶点抽样不确定性。
  \item 本稿为既有文件的回顾性整合，没有重新运行上游模型；这保证了口径唯一，但无法补回从未生成的指标。
\end{enumerate}

\subsection{下一步最小充分验证}

对后续修复模型，建议保持本稿所有阈值和分母不变，先在 17 个复合物每靶点 100 条的盲生成集上冻结序列，再统一生成结构。预先登记的主要终点应是：原始生成分母上的含甲基有效率、全局对齐后环肽姿态 RMSD $<3$ 和 $<5$~\AA 比例；次要终点才是肽自拟合 RMSD、置信度、多样性、通透性预测和同协议天然对照能量。模型修复前后必须使用相同靶点、随机种子规划、结构生成器版本和 RMSD 脚本，且禁止用天然序列恢复率挑选代表。只有这一步通过，才适合扩大到每靶点 1000 条。
